\documentclass[12pt]{article}
\usepackage[T1]{fontenc}
\usepackage[utf8]{inputenc}
\usepackage{lmodern}
\usepackage{textcomp}
\usepackage{enumitem}
\usepackage{geometry}
\usepackage{graphicx}
\usepackage{amsmath, amssymb}
\geometry{margin=1in}
\begin{document}
\begin{center}
Sociology - Undergraduate Level Assessment 21 \
Total Marks: 40 \
Answer all questions.
\end{center}

\section*{Multiple Choice (10 marks)}
\begin{enumerate}[label=\arabic*.]
\item The theory of development which suggests that societies move from traditional to modern, industrial forms of organization is called:
\begin{enumerate}[label=(\alph*)]
    \item Westernization theory
    \item Modernization theory
    \item Industrialization theory
    \item Dependency theory
\end{enumerate}
\item Goldthorpe identified the 'service class' as:
\begin{enumerate}[label=(\alph*)]
    \item those in non-manual occupations, exercising authority on behalf of the state
    \item people working in consultancy firms who were recruited by big businesses
    \item the young men and women employed in domestic service in the nineteenth century
    \item those who had worked in the armed services
\end{enumerate}
\item Role-learning theory suggests that
\begin{enumerate}[label=(\alph*)]
    \item we internalise and take on social roles from a pre-existing framework
    \item we create and negotiate our roles through interaction with others
    \item social roles are not fixed or stable but fluid and pluralistic
    \item roles have to be learned to suppress unconscious motivations
\end{enumerate}
\item Ethnographic research produces qualitative data because:
\begin{enumerate}[label=(\alph*)]
    \item the findings are amenable to statistical analysis
    \item it is conducted over a period of several years
    \item it uncovers rich, detailed accounts from an insider's perspective
    \item it compares findings from a number of different cases
\end{enumerate}
\item 'Scientific' theories in the nineteenth century tried to explain race in naturalistic terms. Which of the following ideas was not considered?
\begin{enumerate}[label=(\alph*)]
    \item genetics
    \item evolution
    \item height
    \item brain size
\end{enumerate}
\end{enumerate}

\section*{True / False (10 marks)}
\textit{Answer True or False}
\begin{enumerate}[label=\arabic*.]
\setcounter{enumi}{5}
\item True or False: Cultural restructuring solely focuses on preserving historical sites without engaging in economic revitalization.
\item True or False: Sociology is a social science because its theories are logical, supported by empirical evidence, and subject to scrutiny by peers.
\item True or False: The 'absolute' poverty line indicates the most extreme level of poverty found in a society.
\item True or False: The post-war welfare state of 1948 did not aim to provide a minimum wage for its citizens.
\item True or False: In the context of the labour movement in the nineteenth century, 'incorporation' involved including working class organizations in political bargaining.
\end{enumerate}

\section*{Long Form Response (20 marks)}
\textit{Answer each question with a clear and complete explanation.}
\begin{enumerate}[label=\arabic*.]
\setcounter{enumi}{10}
\item Discuss how a population pyramid reflects demographic trends, specifically focusing on a population with a sustained higher birth rate than death rate over generations. What shape would you expect, and why?
\item Discuss the limitations of official statistics on strike action in sociology, focusing on the implications of underreported strikes and their effects on social movements.
\end{enumerate}

\vfill
\noindent\textit{End of Paper}
\end{document}